\chapter{«Vint segons»}
\label{sec:1}
Un oficial d'intel·ligència, descrivint el seu paper en un pipeline de selecció d'objectius construït al voltant de llistes de persones per matar generades per una màquina, va explicar a un periodista que dedicava vint segons a cada nom, prou per comprovar que l'objectiu era un home: «No cal gaire temps per saber si algú té veu d'home o de dona.» Va afegir que, com a humà, no aportava cap valor afegit, tret de ser un segell d'aprovació. «Estalviava molt de temps» \autocite{abraham2024lavender}. Segons el seu relat, a això es reduïa l'«humà en el bucle» sota la cadència operativa.\footnote{La recerca sobre el biaix d'automatització prediu aquesta delegació: quan altres exigències competeixen per la seva atenció, les persones que reben la recomanació d'una màquina la comproven menys i s'hi fien més \autocite{parasuraman2010complacency}.}

Una interpretació és que l'oficial és un subjecte a qui es va donar massa poc temps, i que la solució és governar-lo millor: més temps, un segon revisor, una dosi addicional de consciència instal·lada en el bucle. Des d'aquesta òptica de la \emph{subjecció social}, la persona és un subjecte que és fora d'una màquina i la fa servir com una eina: un conductor i un cotxe, un obrer i un torn. La persona té un rol, un nom, responsabilitats i una consciència a la qual es pot adreçar. Una òptica alternativa és la de la \emph{servitud maquínica}, en què la persona és un component de la màquina, acoblat amb les altres peces en un sistema que les regula conjuntament; la màquina fa servir la persona com un dels seus òrgans, i el que la persona aporta és una funció executada al ritme que fixa el sistema.\footnote{Doctorow anomena la posició de l'oficial el centaure invers: no una persona assistida per una màquina, sinó un component d'una màquina fet d'una persona \autocite{doctorow2026reversecentaur}.} Les dues òptiques no s'exclouen; una persona pot tenir un rol en un registre mentre és una peça en l'altre \autocite[Altiplà 13, «The Apparatus of Capture»: hi ha servitud quan els éssers humans són ells mateixos peces constituents d'una màquina que componen entre ells i amb altres coses, sota el control i la direcció d'una unitat superior, i subjecció quan la unitat superior constitueix l'ésser humà com a subjecte vinculat a un nou objecte exterior, que pot ser un animal, una eina o fins i tot una màquina]{deleuze1987plateaus}, i l'oficial dels vint segons ocupava totes dues posicions. Com a peça, era el punt on el pipeline produïa el seu segell, i vint segons eren el temps de cicle d'aquella peça. Com a subjecte, se l'interpel·lava tota l'estona: va descriure a un periodista la dissonància entre les notícies de morts civils i la pressió a la feina per obtenir més objectius. I fins i tot quan la dissonància li arribava a la consciència, va seguir endavant. Es va fer sentir i no va decidir res;\footnote{O, almenys, res que impedís a les Forces de Defensa d'Israel matar 100.000 persones a Gaza entre l'octubre de 2023 i l'octubre de 2025 \autocite{gomezugarte2025uncertainty,acosta2026gaza,economist2025gaza}. I el relat d'aquell oficial en concret el va publicar +972 Magazine l'abril de 2024. Ningú no sap si la decisió d'aquell oficial de parlar amb un periodista va evitar una mort, 50.000 morts o cap. I què hauria passat si l'oficial s'hagués negat després de dues setmanes en el càrrec, posem per cas, hipotèticament, el novembre de 2023? Parlem d'un oficial d'intel·ligència que feia servir Lavender per elaborar llistes de persones per matar. El seu entorn de treball quotidià probablement s'assemblava més a una feina d'oficina civil que a una pel·lícula d'acció. Hauria d'haver demanat reunions i més reunions amb els seus superiors, amunt per la cadena de comandament, alhora que es negava a aprovar objectius per a atacs aeris dels quals es podia esperar raonablement que causessin morts civils? Es pot defensar que això hauria estat el que era ètic fer. En aquest supòsit contrafactual, hi hauria hagut menys morts?} el judici que la política era desproporcionada li va arribar, va dir, en retrospectiva.

Així doncs, el que faltava no era una consciència; en tenia una. Una consciència funciona fent que una decisió costi alguna cosa, però aquí els costos anaven en sentit contrari: «Ens pressionaven constantment: “Porteu-nos més objectius.” De debò que ens cridaven» \autocite{abraham2024lavender}. Era un component amb escrúpols, però la cadència es mantenia: els escrúpols no canviaven res del que aportava, i es va convertir en un segell. Un component amb escrúpols que no manté la cadència es converteix en el «coll d'ampolla humà» que un antic comandant de la unitat d'intel·ligència de senyals d'Israel va dedicar tot un llibre a eliminar \autocite{ys2021humanmachine}. El que va fer de l'oficial un component va ser el dispositiu, no la seva actitud. Un component pot seguir endavant a contracor o amb entusiasme: el funcionari irreflexiu que Arendt va veure al judici d'Eichmann \autocite{arendt1963eichmann} era, a jutjar pel que havia dit a l'Argentina quatre anys abans, quan no s'hi jugava res, un paper que interpretava per a Jerusalem \autocite{stangneth2014eichmann}. El que el pipeline deixava aportar a l'oficial era una operació fixa a un ritme fix: veia cada cas tal com el descrivia el pipeline, un per un, en el temps que el pipeline li donava. Si alguna cosa podria ocupar aquesta posició sense ser un component és la pregunta de la part IV.

Diverses investigacions periodístiques van descriure tres sistemes que funcionaven conjuntament a Gaza: un que generava objectius entre edificis i infraestructures, i segons les fonts la majoria dels suggeriments eren habitatges particulars; un segon que puntuava individus segons la seva presumpta afiliació a un grup armat i produïa llistes que, pel que sembla, van arribar a marcar fins a 37.000 persones, i que internament se sabia que s'equivocaven aproximadament una vegada de cada deu\footnote{L'exèrcit israelià va contradir aquest relat. Els relats de Minab i de Maven que segueixen també es basen en periodisme d'investigació que les parts igualment discuteixen. Els prenc tal com s'han publicat.}; i un tercer que seguia quan les persones marcades tornaven a casa, perquè els atacs aeris es poguessin fer coincidir amb els moments en què eren amb la família.

\runin{Maven: la versió nord-americana, ara accelerada amb IA generativa}

El bombardeig de l'escola de Minab va matar 156 persones, 120 de les quals eren escolars \autocite[par.~92]{UNFFMIran2026}.\footnote{Les xifres varien segons la font i la data, de 156 (el balanç definitiu del fiscal de Minab, que també dona Amnistia Internacional) a 168 (l'alcalde de la ciutat, els primers dies). La missió de l'ONU dona les xifres del fiscal i també recull informació independent de 157 morts, 123 dels quals infants de 13 anys o menys \autocite[par.~92–93]{UNFFMIran2026}.} La funció civil de l'escola era visible en imatges de satèl·lit de fonts obertes, i la codificació de l'objectiu encara duia una designació heretada d'una instal·lació contigua de la Guàrdia Revolucionària. Una missió de determinació dels fets de l'ONU va informar que hi havia motius raonables per creure que el bombardeig constituïa el crim de guerra de fer d'un bé de caràcter civil l'objecte d'un atac \autocite{ohchr2026minabpress}, i que l'escola era clarament identificable, sense cap informació que n'indiqués un ús militar \autocite[par.~81–86, 124–126]{UNFFMIran2026}. Va concloure que els Estats Units van incomplir l'obligació de fer tot el que era factible per verificar l'objectiu, va comptar l'actualització de les dades d'objectius entre les precaucions factibles i va situar l'incompliment més enllà de la negligència \autocite[par.~124–125]{UNFFMIran2026}. Amnistia Internacional va arribar a una conclusió paral·lela, en uns termes que segueixen el deure de verificació de l'article 57(2)(a)(i) del Protocol addicional I: que les forces potser s'havien basat en informació d'intel·ligència obsoleta i no havien fet tot el que era factible per verificar l'objectiu \autocite{Amnesty2026Minab}.

Va passar el matí del 28 de febrer, el primer dia de la campanya de l'Iran \autocite{Amnesty2026Minab,crs2026iranwar}. El dia abans, el secretari havia anunciat que el proveïdor (Anthropic) el model del qual (Claude) era dins la plataforma de selecció d'objectius (Maven Smart System) seria designat un risc per a la cadena de subministrament per haver-se negat a concedir llicència de Claude per a qualsevol ús legal \autocite{npr2026openai}. Maven Smart System havia proposat l'escola com a objectiu recomanat per al primer dia \autocite{livingstone2026maven,gizmodo2026minab}. Cap relat publicat no diu quin paper va tenir Claude en aquesta recomanació, i una persona coneixedora de la feina del proveïdor va dir a NBC News que el model no donava directament consells de selecció d'objectius \autocite{nbc2026targeting}.

Una revisió interna del Pentàgon, les conclusions completes de la qual es mantenen en estricta reserva, hauria constatat que la preparació d'objectius que en conflictes anteriors duia hores s'havia condensat en minuts, que el personal de mitigació de danys civils havia caigut aproximadament un noranta per cent, fins a menys de vint persones, que l'equip del Comandament Central havia passat de deu persones a una, i que cap dels seus membres no havia revisat l'emplaçament de Minab, un resultat que els responsables van atribuir tant a les retallades com a la decisió de no implicar l'equip en la planificació i les operacions \autocite{Bloomberg2026KillChain}. La investigació formal s'havia obert al març, després que una avaluació preliminar conclogués que els Estats Units n'eren probablement responsables \autocite{NPR2026MinabProbe}.

A la campanya de l'Iran del 2026, una plataforma de selecció d'objectius construïda per un contractista, amb el model d'ús general d'un proveïdor comercial a dins, que es feia servir per classificar objectius per importància estratègica i redactar una justificació jurídica per a cada atac \autocite{livingstone2026maven}, va ajudar a atacar més de mil objectius en les primeres vint-i-quatre hores, unes deu vegades el que era possible abans \autocite{csis2026maven}, i tretze mil al dia trenta-vuit \autocite{breakingdefense2026maven}. El Pentàgon no n'ha publicat cap temps de revisió per objectiu.\footnote{Maven és un sol nom de programa que porta més d'un sistema: el programa del Pentàgon que encara es diu Maven, i el Maven Smart System, la plataforma de programari que en va sortir \autocite{csis2026maven}. El programa va començar el 2017 com un projecte per etiquetar objectes als vídeos que els drons recollien durant la campanya contra l'Estat Islàmic \autocite{weisgerber2017maven,fang2019maven}. El sistema del 2026 fusiona sistemes militars i fonts d'intel·ligència en una única interfície que serveix per trobar i atacar objectius i per comandar operacions \autocite{defensescoop2026transition}. La plataforma és de Palantir: el cap digital i d'IA del Pentàgon l'anomena «el producte de programari com a servei de Palantir» \autocite{defensescoop2026transition}. La posició pública de Palantir, citada de segona mà, és que decidir sobre la vida o la mort no és el seu paper i que la responsabilitat sempre és de l'organització militar \autocite{lee2026palantir}, i l'empresa sosté que la intel·ligència de base era responsabilitat del govern i que no hi ha cap prova que Maven funcionés malament ni que causés el bombardeig de Minab \autocite{Bloomberg2026KillChain}. Com va passar el programa d'una oficina i un contractista al següent s'explica a l'apèndix~\ref{sec:C.1}.}

Les persones de les llistes d'atacs aeris israelianes i nord-americanes no haurien pogut objectar que hi fossin: els pipelines no estaven construïts per recollir aquesta mena de resposta. Andry José Hernández Romero va quedar atrapat en un altre pipeline, i ell sí que hauria pogut objectar. Maquillador veneçolà, es va presentar al punt d'entrada de San Ysidro l'agost de 2024 per demanar asil, va superar l'entrevista de por creïble i tenia advocat, una sol·licitud pendent i data de vista \autocite{LULAC2025HernandezRomero,Vega2025HernandezRomero}. L'empleat d'un contractista que el va examinar va llegir dos tatuatges, unes corones damunt les paraules «mama» i «papa», com a marques d'una banda \autocite{Palomera2025Blade}, i el març de 2025 el van enviar en avió a una presó del Salvador abans de la seva vista, el mateix vespre que un jutge federal ordenava que els avions fessin mitja volta \autocite{jgg2025contempt}. Tenia tot el que una persona necessita per objectar, excepte el temps. Tornarem al seu cas més a fons quan parlem del pipeline de segrest i trasllat forçós dels Estats Units.

\textbf{Qui es va negar.} L'expedient dels vint segons té una segona meitat, l'expedient de qui va resistir pressions d'aquesta mena, i diu que la composició de la sala importa molt. En la condició de base de Milgram, aproximadament un terç dels subjectes es van aturar abans de la descàrrega màxima. Amb dos companys que desobeïen, només un de cada deu va arribar fins al final, i quan dos experimentadors discrepaven, en van continuar molts menys que en la condició de base \autocite{milgram1974obedience}. Aproximadament una quarta part dels subjectes d'Asch no van cedir mai davant la majoria equivocada \autocite{asch1956studies}. Browning calcula que entre una desena i una cinquena part dels homes del Batalló 101 de la Policia de Reserva es van negar a participar en els afusellaments de Józefów, la primera massacre del batalló, o se'n van escapolir \autocite{browning1992ordinary}. La proporció anava de gairebé ningú a gairebé tothom segons l'entorn. Les persones que van resistir tenien proves de fora del canal: un company que desobeïa, una segona autoritat que discrepava i, entre els rescatadors que van estudiar els Oliner, relacions que travessaven la línia que traçava la persecució \autocite{oliner1988altruistic}. O bé mantenien el compromís amb una precisió que la pressió no podia vèncer, tal com es mantenen els mandats morals i els valors sagrats \autocite{skitka2009limits,ginges2007sacred}. O bé consideraven que la font hi tenia interessos. I els grups resistien allà on els individus aïllats rarament ho feien \autocite{gamson1982encounters}.

La sala també es pot compondre a favor de l'altre bàndol. L'Alemanya de Weimar tenia una vida associativa densa, que probablement va ajudar a fer caure la república: a mesura que els partits i els governs no responien als greuges de la gent, la gent posava l'energia en clubs i societats organitzats per grups, i els nazis van reclutar a través de les xarxes que aquestes societats ja havien construït \autocite{berman1997weimar}. Les ciutats amb més corals, clubs esportius i associacions gimnàstiques van veure una entrada més ràpida al partit, però només allà on el govern era inestable; a Prússia, que tenia un govern estable, el vincle era petit i no significatiu \autocite{satyanath2017bowling}. Així doncs, el que tenien els rescatadors no era simplement més companyia. Els seus vincles travessaven la línia que traçava la persecució. Ashutosh Varshney va trobar el mateix a les ciutats índies: els aldarulls entre hindús i musulmans eren rars allà on les associacions aplegaven les dues comunitats en les mateixes organitzacions, i freqüents allà on no ho feien \autocite{varshney2002ethnic}. Els vincles que fan possible la negativa són els que un moviment feixista hauria de tallar.

Tot el que segueix tracta de com compondre la sala.\footnote{La sala és l'escenari, i com diria el meu professor de teatre, l'escenografia dirigeix l'escena. Qui més hi ha a la sala? Es pot marxar i sortir-ne viu? Algú sentiria una negativa?} Compondre és una paraula que vol dir fer, i Hannah Arendt advertia de no tractar la política com una cosa feta. Per als grecs, escrivia, la llei era literalment un mur, sense el qual hi podia haver una població però no una ciutat \autocite{arendt1958human}: el legislador construïa el mur, i el que passava a dins no li tocava a ell fer-ho. Les mesures bàsiques que defenso són murs en aquest sentit. Permeten negar-se i sortir-ne viu, i fan que la negativa se senti; no la produeixen, i el poder, segons ella, només existeix entre persones que actuen juntes i desapareix quan es dispersen, que és el que l'expedient de qui es va negar mostra sobre la negativa. Els instruments de les parts II i III fan que una negativa compti. Les institucions de la part V fan insubmisos. Una màquina que es pogués negar necessitaria el que tenien els nou de cada deu amb un company al costat: un company amb una història diferent al seu costat, un lloc on anar i algú que la sentís.

Això és també el que un insubmís —incloent-hi una màquina insubmisa— podria ser per a l'oficial, dins el pipeline. Una màquina que es negués seria a la pantalla quan el cas li arribés: el company que desobeeix, si l'efecte sobreviu al pas d'una sala a una pantalla, cosa que no s'ha provat mai i que el biaix d'automatització prediu que no passarà. Una màquina que mogués els revisors per deferència n'hauria fet segells d'un altre signe. L'oficial amb vint segons i una màquina que es nega al seu costat és el mateix oficial que el del reportatge?
